\section{Introduction}
\label{sec:introduction}

Escalating capabilities within a society create escalating
mechanisms for disaster the civilisation has never experienced,
some existentially catastrophic if blindly accelerated toward.
This capability-escalating regime is precisely the
$\volL$-maximising dynamic that Goal-Frontier Maximization (GFM)
seeks to produce; the question this paper addresses is how to
balance the existential risks of such escalation against the
reciprocal danger of a
society that becomes increasingly bound by restrictive
overcaution, accumulating restrictions it cannot revise through
observation, approaching the domination feedback loop
of~\citep[Loop~3]{lasser2026phase}.

The classical balancing mechanism is catastrophe-as-teacher: a
first occurrence writes historical markers---tsunami stones,
disaster records, institutional memory, building-code revisions,
pharmacovigilance regimes---that subsequent generations consult in
lieu of the experiment they cannot run.  This mechanism is
unavailable for capabilities novel to any generation, which is
precisely the regime the escalating dynamics generate.  This paper
constructs a formal alternative: a mechanism for generating the
evidential content historical markers would have recorded, within
a bounded damage budget chosen by the framework rather than the
next generation.

The specific failure mode in the existing GFM framework is the
max-aggregation pathology of~\citep{lasser2026horizon}, named
there as the most important open problem of the sequence: the
conservative max operator, combined with the verification
asymmetry inherent in risk prevention, produces overcaution that
is structurally irresolvable within the existing machinery---a
risk claim that successfully prevents catastrophe generates no
observation that could lower the risk estimate, because the very
success of the restriction is what prevents the counterfactual
test.  This is a distinct problem from the compound feedback
loops analysed in~\citep{lasser2026phase}: that work asks when
partial subsumption self-corrects versus cascades, while the
present paper asks a complementary question---when the framework
identifies a risky capability and imposes a restriction, how
does it learn whether the restriction remains necessary, or
whether technological progress has mitigated the underlying
risk?  The two problems reinforce each
other:~\citep[Proposition~4]{lasser2026phase} shows the
restriction-verification gap is a degenerate phase boundary with
zero self-correction rate ($\rS = 0$), but the phase-boundary
machinery does not by itself supply the missing self-correction
mechanism.

The paper is named for Mayor Kotoku Wamura's tsunami-preparedness
legacy at Fudai~\citep[Discussion]{lasser2026poset}: an oversized
seawall built on the basis of centuries-old stone markers warning
against construction below a certain elevation, ridiculed for
decades as wasteful overcaution, and vindicated when the 2011 Tohoku
tsunami reached precisely the elevation the markers had recorded.
The Fudai case describes the shape of the problem; the question
this paper answers is how to have Wamura's foresight without the
ancient guidance that informed him.

\subsection{The structural problem: long-timeline risks}

The ratchet arises from a timeline mismatch between the risk and
the framework's observational window.  Let $T$ be the realisation
timeline of a risk---the interval on which the risk would actually
manifest under exercise---and $W$ be the decision-cycle
observational window over which the framework can accumulate
evidence.  When $T \ll W$, observation resolves the question
promptly; the restriction dynamics are well-behaved.  When
$T \gg W$, no observation can update the risk estimate within any
practical decision cycle.  The restriction is adopted under
uncertainty, survives because no counter-evidence can arrive in
time to challenge it, and restrictions of this class accumulate
monotonically because: (i) new long-timeline risks are continually
discovered, each imposing a restriction; (ii) existing long-timeline
restrictions cannot be revised through observation.  The ratio of
restriction-imposing decisions to restriction-revising decisions
approaches infinity on the long-timeline subset.

Three mechanisms compose with the timeline mismatch to produce the
pathology.

\paragraph{Max-aggregation.}
The risk framework of~\citep{lasser2026horizon} takes the maximum
over agents' risk estimates, so a single high-risk claim from any
agent constrains the capability tuple.  This is conservative by
design---it prevents the framework from dismissing minority risk
warnings---but it means a single overly cautious agent can lock a
capability indefinitely when the underlying risk is long-timeline.

\paragraph{Verification asymmetry.}
The monotonicity
result~\citep[Proposition~1]{lasser2026exo} shows that as an
agent's capabilities grow, the verification asymmetry between what
it can demonstrate and what can be externally verified is
non-decreasing.  Applied to risk claims: verifying a risk (observing
catastrophe) is operationally straightforward on any timeline;
verifying safety (observing non-catastrophe under exercise) requires
the exercise to happen, which the restriction prevents, and even
when permitted would not produce a conclusion within $W$ when
$T \gg W$.

\paragraph{Neutral-prior bootstrap.}
When a new capability is discovered and no exercise history exists,
the framework initializes with a neutral
prior~\citep[Discussion]{lasser2026horizon}.  Combined with
max-aggregation, a single agent assigning high risk to the new
capability immediately restricts it, and the restriction then
prevents the evidence accumulation that would revise the estimate.

\medskip

The three mechanisms compose into a ratchet: risk claims can raise
$\pR$ but the restricted regime provides no mechanism to lower it
within the framework's operating horizon.
We formalize this ratchet in
Section~\ref{sec:verification_asymmetry}, introduce controlled
relaxation protocols as the structural resolution in
Section~\ref{sec:relaxation_protocol}, and prove the damage-bound
and convergence theorems in
Sections~\ref{sec:damage_guarantees}--\ref{sec:convergence}.

\subsection{Framing precondition: $\volL$ as an operational target}
\label{sec:framing_precondition}

All results in this paper are stated in terms of $\volL$ dynamics
and Bayesian updates on risk claims about $\volL$-contraction.  We
assume throughout that $\volL$ is an operationally meaningful
target, in the sense inherited
from~\citep{lasser2026gfm,lasser2026phase}: improving $\volL$
corresponds to an improvement in the experiential capability volume
the framework is designed to protect.  This is the precondition
under which the damage-bound and convergence claims carry their
safety interpretation.  If $\volL$ diverges from the experiential
target (the B-to-C gap~\citep{lasser2026gfm}), the damage tolerance
$\damagebound$ bounds contraction on the wrong quantity and every
subsequent result inherits the divergence.  Controlled relaxation
does not close the B-to-C gap; it resolves the restriction-to-risk
feedback gap that exists \emph{above} the B-to-C precondition.

\subsection{Summary of results}

\begin{enumerate}
\item \textbf{Formalization of the Wamura asymmetry}
  (Section~\ref{sec:verification_asymmetry}).
  We define risk-restriction pairs and their two observation
  regimes, then prove that under the restricted regime the
  Bayesian update has trivial likelihood ratio: successful
  prevention is observationally indistinguishable from unnecessary
  prevention (Proposition~\ref{prop:asym_evidence}).  Combined
  with max-aggregation and neutral-prior initialization, this
  produces monotone overcaution
  (Proposition~\ref{prop:monotone_overcaution}) and an asymptotic
  lockup corollary.

\item \textbf{Controlled relaxation protocol}
  (Section~\ref{sec:relaxation_protocol}).
  A controlled relaxation test lifts the restriction within a
  bounded sub-poset $\scope$ for a fixed duration $\testdur$,
  under continuous multi-channel monitoring, with an
  instantaneous-damage safety cutoff at the tolerance
  $\damagebound$.  The test is formalized as a simple binary
  hypothesis problem ($H_0$: pathway harmless, $H_1$: pathway
  dangerous with fixed activation probability $\thetaone$) with
  fixed per-test likelihood ratios $\Lzero < 1 < \Lone$.

\item \textbf{Damage bound}
  (Section~\ref{sec:damage_guarantees}).
  Theorem~\ref{thm:damage_bound} bounds the worst-case $\volL$
  contraction during a test by $\damagebound$ during the test
  window and $2\damagebound$ post-reinstatement, under a scope
  construction (Definition~\ref{def:test_scope}) whose
  cooperative-closure condition explicitly accounts for
  cross-boundary forward-reachability.

\item \textbf{Convergence theorem}
  (Section~\ref{sec:convergence}).
  Theorem~\ref{thm:convergence} shows that under repeated
  controlled relaxation with fixed parameters, the posterior
  $\pR$ converges almost surely to the correct hypothesis, with
  exponential rate given by the Kullback--Leibler divergence
  between the observation distributions under $H_0$ and $H_1$.
  Finite-sample corollaries give explicit bounds on the number of
  tests required to clear a false risk claim.

\item \textbf{Integration}
  (Section~\ref{sec:integration}).
  The protocol preserves the self-balancing property of $\volL$
  (Proposition~\ref{prop:preserve_balance}), provides empirical
  arbitration for inter-agent risk disagreement
  (Proposition~\ref{prop:empirical_arbitration}), and extends the
  multi-channel detection architecture
  of~\citep{lasser2026scm} with a sixth, actively generated
  channel (Lemma~\ref{lem:sixth_channel}).
\end{enumerate}

\subsection{Real-world interpretation}
\label{sec:real_world}

The formal construction of this paper is a programmatic mechanism
for empirical review of long-timeline restrictions that may have
outlived their justifying risk model.  The framework identifies
restrictions that persist without recent confirming evidence,
formulates the underlying risk pathway that the restriction was
intended to block, and designs a bounded-damage empirical test that
probes whether the physical mechanism underlying the restriction
still operates---a test conducted by appropriately trusted actors
within a bounded scope, with the trust apparatus
(cryptographic~\citep{lasser2026exo}, governmental, institutional,
legal, or otherwise) selected to match the domain.

Application domains where long-timeline restrictions accumulate
and benefit from programmatic review:

\begin{itemize}
\item \textbf{Environmental contamination controls}: chemical
  scheduling where bioaccumulation or persistence timelines are
  multi-decade (DDT, PFAS, ozone-depleting refrigerants); the
  current restriction may target a mechanism that has changed
  under reformulation or that modern mitigation has obsoleted.
\item \textbf{Nuclear dual-use controls}: export and technology
  controls where weapons-grade and reactor-grade pathways share
  keystone capabilities in broad restriction regimes designed
  under mid-20th-century technology.  Modern enrichment and
  machining have differentiated to a degree that supports
  narrower keystone choices; see
  Section~\ref{sec:keystone_refinement}.
\item \textbf{Pandemic preparedness and dual-use biology}:
  research moratoria on gain-of-function work, where the risk
  timeline is the interval between pandemic events (decades)
  and the observational window is research-cycle short.
\item \textbf{Infrastructure safety margins}: dam, bridge, and
  seismic design codes originally calibrated against
  century-scale events, where construction material and
  mitigation-technology progress may permit refined margins.
\item \textbf{AI capability restrictions}: the sequence's
  primary application.  Restrictions adopted against
  concerns about future outcomes when no empirical
  counter-evidence can arrive within training and deployment
  cycles.
\end{itemize}

This list is illustrative; the formal content does not commit to
any particular domain.  The framework is indifferent to whether
the trusted actors are human oversight bodies, cryptographically
committed automated systems, or combinations thereof; the
implementation of trust is a domain-specific design decision
outside the scope of the formal construction.

\subsection{What this paper does not do}
\label{sec:not_do}

\begin{itemize}
\item
\textbf{It does not replace max-aggregation.}  Log-odds risk
aggregation~\citep[Definition~5]{lasser2026scm} already provides
an alternative for the multi-agent case.  The verification
asymmetry we resolve persists under \emph{any} aggregation rule
whenever a risk claim prevents the exercise that would generate
evidence.  Controlled relaxation is orthogonal to the choice of
aggregator.

\item
\textbf{It does not address observation-channel integrity under
adversarial conditions.}  The commitment protocols
of~\citep{lasser2026exo} handle that.  We assume honest
observation channels and honest (but possibly miscalibrated) risk
assessors throughout.

\item
\textbf{It does not resolve the B-to-C gap.}  Controlled
relaxation generates evidence about $\volL$-contraction claims,
not about experiential value.  Convergence of $\pR$ to truth
translates into convergence of the restriction landscape only
insofar as $\volL$ tracks the experiential target; the
precondition from Section~\ref{sec:framing_precondition} is
inherited.

\item
\textbf{It does not authorise unconstrained exercise of dangerous
capabilities.}  A controlled relaxation test occurs strictly
within a bounded scope $\scope$, under continuous detector
monitoring, with an instantaneous-damage cutoff at the user-set
tolerance $\damagebound$---a clinical-trial protocol, a contained
laboratory, a sandboxed deployment, an instrumented test range,
or other domain-appropriate constrained environment.  The
framework does not propose detonating bombs occasionally to
confirm bombs still explode: risks with abundant current evidence
are screened out by the meta-policy of
Section~\ref{sec:meta_policy} before any test is initiated,
because the information value of confirming an established
mechanism is near zero while the expected damage remains positive.
The protocol addresses the complementary case, where evidence is
absent and restrictions accumulate without challenge.

\item
\textbf{It does not claim that controlled relaxation is always
preferable to structural resolution.}
Section~\ref{sec:meta_policy} places the test-initiation decision
against the information value of a structural
argument~\citep[Proposition~2]{lasser2026horizon}: when a
structural resolution is available (for example, a proof that
the pathway is mechanically infeasible), it is strictly
preferable to any controlled relaxation test.

\item
\textbf{It does not address the Fudai case directly.}  At Fudai,
historical evidence (ancient tsunami stones) existed but was
dismissed in favour of the modern observational window.  The
remedy there is to restore historical records to the decision
process, not to generate new experimental evidence---controlled
relaxation would have been actively dangerous at Fudai (a test
of a reduced seawall works perfectly until the tsunami arrives).
The paper's domain is the complementary case where no historical
record is available.
\end{itemize}

\subsection{Dependencies on prior results}

Table~\ref{tab:dependencies7} summarizes which results
from~\citep{lasser2026gfm, lasser2026poset, lasser2026horizon,
lasser2026scm, lasser2026exo, lasser2026phase} this paper relies
on and how they enter the analysis.

\begin{table}[h]
\centering
\small
\begin{tabular}{@{}p{2.0cm}p{5.6cm}p{6.0cm}@{}}
\toprule
\textbf{Source} & \textbf{Result used} & \textbf{Role here} \\
\midrule
Poset        & Axiom verification (Prop.~1), benchmark refinement
             & $\volL$ properties preserved under relaxation;
               complementary evidence source \\
Horizon      & Exercise Pathway (Def.~2), Concentration Risk (Def.~3),
               Risk Claim (Def.~4), Preemptive Restriction (Prop.~1)
             & The risk machinery the protocol operates within \\
Horizon      & Structural discovery value (Prop.~2), anti-monopolar
               property (Prop.~6)
             & Value framework for test-cost vs.\ restriction-cost
               comparison \\
Causal       & Risk-Trust Dynamics (Def.~4), Log-Odds Risk
               Aggregation (Def.~5), $L^2$ Convergence (Prop.~2)
             & Bayesian update machinery that processes test
               outcomes \\
Causal       & Multi-channel architecture (Table~1), Causal
               Attribution (Def.~2)
             & Formalizing what ``observing the outcome'' means
               during a test \\
Verification & Verification Asymmetry (Def.~2), Monotonicity
               (Prop.~1), Exercise Protocol (Def.~12), Risk Claim
               Protocol (Def.~13)
             & The formal verification-asymmetry framework the
               protocol resolves \\
Verification & Channel Isolation (Prop.~5), Channel Independence
               (Prop.~6)
             & Integrity conditions on test observations \\
Phase        & Phase boundary (Thm.~1), Wamura problem as
               degeneracy (Prop.~4)
             & Controlled relaxation makes $\rS > 0$ for Loop~3 \\
\bottomrule
\end{tabular}
\caption{Dependencies on prior results.
  Poset~=~\citet{lasser2026poset},
  Horizon~=~\citet{lasser2026horizon},
  Causal~=~\citet{lasser2026scm},
  Verification~=~\citet{lasser2026exo},
  Phase~=~\citet{lasser2026phase}.}
\label{tab:dependencies7}
\end{table}
